\chapter{Analyse des besoins, sp\'ecifications et m\'ethodologie}

\guided{Le lecteur doit pouvoir reconstituer la cha\^ine logique besoin $\rightarrow$ exigences $\rightarrow$ m\'ethode. Ce chapitre est essentiel pour montrer la ma\^itrise du cadrage d'ing\'enierie.}

\section{Analyse de l'existant et du contexte d'usage}
\todo{Pr\'esenter l'environnement r\'eel du projet, les utilisateurs, les contraintes organisationnelles, techniques ou r\'eglementaires.}

\section{Besoins fonctionnels et non fonctionnels}
\begin{table}[H]
\centering
\caption{Exemple de formalisation des besoins}
\begin{tabularx}{\textwidth}{>{\RaggedRight\arraybackslash}p{1.8cm} >{\RaggedRight\arraybackslash}X >{\RaggedRight\arraybackslash}p{3cm} >{\RaggedRight\arraybackslash}p{2.2cm}}
\toprule
ID & Besoin / exigence & Type & Crit\`ere d'acceptation \\
\midrule
EX-01 & \todo{Mesurer / d\'etecter / pr\'edire ...} & Fonctionnelle & \todo{Pr\'ecision attendue} \\
EX-02 & \todo{Temps de r\'eponse} & Non fonctionnelle & \todo{$< 200$ ms} \\
EX-03 & \todo{S\'ecurit\'e / robustesse / fiabilit\'e} & Contrainte & \todo{Condition de validation} \\
\bottomrule
\end{tabularx}
\end{table}

\section{Sp\'ecifications et indicateurs de performance}
\todo{Relier les exigences \`a des m\'etriques ou des tests futurs.}

\section{M\'ethodologie retenue}
\todo{D\'ecrire la d\'emarche : cycle de d\'eveloppement, protocole exp\'erimental, outils, jalons, jeux de donn\'ees, banc d'essai, mod\`ele de conception, etc.}

\subsection{Orientation disciplinaire}
\guided{\disciplinehint}
\todo{Ajouter ici les diagrammes ou mod\`eles adapt\'es \`a la fili\`ere : UML/SysML, sch\'ema fonctionnel, architecture r\'eseau, diagramme de flux de donn\'ees, FAST, SADT, diagrammes d'\'etats, etc.}
